% compendium/execution-trading-costs/trade-signing.tex
% Entry migrated on 2026-10-16 from the legacy trade-classification-and-signed-flow file (slice 33), owned by
% price-impact-execution, one of the file's map candidates. Dispositions in notes/migration-slices.md; the file
% verbatim in notes/artefacts/price-impact-execution/.

\section{Signing trades}

\begin{qfmethod}[Trade-sign classification: the quote rule, the tick rule and Lee--Ready, block prices, signed volume, and what an inferred sign does and does not say]\label{meth:trade-sign-classification}
\qffield{Statement} Trade data rarely say which side initiated a trade, so the sign $q_t=+1$ (buyer-initiated) or $-1$ is inferred. If all executions occur at the posted quotes, $q_t=+1$ at the ask and $-1$ at the bid; more generally $q_t=\operatorname{Sign}(p_t-m_t)$ with $m_t$ the prevailing midpoint (the \emph{quote rule}), which leaves midpoint trades unclassified; for those, the trade takes the sign of the most recent price change (the \emph{tick rule}) \citep[chap.~9, sec.~9.7.1, p.~89]{s6-hasbrouck}. Harris describes the combination as the Lee and Ready algorithm: above the midpoint a buyer was aggressive, below it a seller; at the midpoint, an uptick or zero uptick means a buyer, otherwise a seller \citep[chap.~21, p.~423]{alpha-harris}. For block trades, a print closer to the bid than to the ask was probably seller-initiated, otherwise buyer-initiated, presuming the initiator paid a concession \citep[chap.~15, sec.~15.5, p.~332]{alpha-harris}. Signed volume $q_tV_t$, or a concave transform such as $q_t\sqrt{V_t}$, enters impact models; Hasbrouck advises against convex transforms such as $q_tV_t^2$, which amplify extreme volumes (sec.~9.7.2, p.~89).

\qffield{Conditions} Signing needs the quotes and trades in the right order: in limit-order markets the interleaving is usually accurate, but with hybrid markets, different dissemination paths, coarse time stamps or manual steps it is ``likely to be suspect'' (Hasbrouck, p.~89). The rule at the posted quotes works only where the book can be rebuilt and executions are against visible orders, not hidden or reserve ones (p.~89). Lee and Ready (1991) assess the quote and tick tests; Hasbrouck cites them (note~1 to chap.~9, p.~180), and the paper itself was not read here.

\qffield{Practice} Checked (\nolinkurl{order-flow-checks.py}, ten assertions on prices held as exact decimals): midpoint trades after an uptick, or after a zero tick whose last non-zero change was up, are buys, and after a downtick or a zero downtick sells, for whole-tick midpoints (10.01 in 10.00--10.02) and half-tick ones (10.005 in 10.00--10.01); a midpoint trade with no earlier price change is left unclassified; trades just above or below a half-tick midpoint (10.006, 10.004) follow the quote rule whatever the tick; a trade at the bid is a sell, and at the ask a buy, whatever the previous tick. An implementation must compare prices exactly: in binary floating point $(10.00+10.01)/2$ is 10.004999\dots, below 10.005, so the half-tick midpoint trade after a downtick is wrongly called a buy, as the slice-33 version of the check did. Sign persistence in event time is a stylised fact (\hyperref[ev:high-frequency-stylised-facts]{stylised facts}); Hasbrouck notes that signed size often adds little explanatory power beyond the sign itself (sec.~9.7.2, p.~90). Summed over an interval, signed volume is the trade imbalance of \hyperref[def:order-flow-imbalance]{order-flow imbalance}, not the event-based OFI.

\qffield{Pitfalls} The sign identifies which side demanded immediacy at that moment, not who was better informed or more eager: a buyer who pegs the bid can accumulate aggressively while the algorithm calls every trade passive (Harris's sidebar, p.~423). Measured as transaction costs, the algorithm overstates the costs of traders who do not only take liquidity, and cannot see that several trades filled one order, so it understates that order's impact cost (p.~423). Harris reports that ``analysts have determined that sellers initiate approximately 80 percent of all large block trades in the U.S.\ stock markets'', and reads the asymmetry as consistent with three of his four block-trading problems: latent demand (buyers can buy only from holders when short selling is hard), price discrimination (sellers can credibly reveal their full size) and asymmetric information (large buyers are presumed informed) (pp.~332--333); he gives no sample or period, and the primary evidence is a backlog item. The legacy taxonomy listed the tick rule without its use at the midpoint, and its ``false persistence'' from misclassification and order splitting was not found in these passages.

\qffield{See also} \hyperref[def:spread-measures]{spread measures} (effective spreads need signs); \hyperref[meth:information-driven-bars]{information-driven bars} (the tick rule in bar construction); \hyperref[meth:order-flow-impact-regression]{order-flow impact regression}; \hyperref[mod:propagator]{propagator model}; \hyperref[def:block-trade]{block trade}.

\qffield{Further reading} Lee and Ready (1991) and the studies behind Harris's 80 percent are not read. The legacy file is kept verbatim in \nolinkurl{notes/artefacts/price-impact-execution/}.

\qffield{Index} trade sign; trade classification; buyer-initiated; seller-initiated; quote rule; quote test; tick rule; tick test; zero uptick; Lee--Ready algorithm; midpoint trades; block-price rule; signed volume; trade imbalance; aggressive and passive traders; block initiator asymmetry

\qffield{Status} Verification: partly checked 2026-10-16 (against the text layers of local copies of \texttt{s6-hasbrouck}, 2007 Oxford printing, pp.~89--90 and 180, and \texttt{alpha-harris}, 2002 Oxford printing, pp.~332--333 and~423; the rule computed on toy trades, recomputed 2026-10-16 with exact decimal prices after a floating-point error in the first check). Stability: draft. Currency: methods timeless; the block statistic as reported in 2002.
\end{qfmethod}
